\section{Теоретические основы}

\subsection{Аналитическое решение для однородного поля}

Рассмотрим движение материальной точки массой \(m\) в однородном гравитационном поле. Пренебрегая сопротивлением воздуха и другими силами, запишем второй закон Ньютона:

\begin{equation}
    m\ddot{\vec{r}} = m\vec{g}
    \label{eq:newton_law}
\end{equation}

где \(\vec{g} = (0, 0, -g)\) — вектор ускорения свободного падения, \(g \approx 9.81 \, \text{м/с}^2\).

В проекциях на оси координат система уравнений принимает вид:

\begin{equation}
\begin{cases}
    \ddot{x} = 0 \\
    \ddot{y} = 0 \\
    \ddot{z} = -g
\end{cases}
\label{eq:projections}
\end{equation}

Интегрируя уравнения \eqref{eq:projections} с учетом начальных условий \(\vec{r}(0) = (x_0, y_0, z_0)\) и \(\vec{v}(0) = (v_{0x}, v_{0y}, v_{0z})\), получим кинематические уравнения движения:

\begin{equation}
\begin{cases}
    x(t) = x_0 + v_{0x} t \\
    y(t) = y_0 + v_{0y} t \\
    z(t) = z_0 + v_{0z} t - \dfrac{g t^2}{2}
\end{cases}
\label{eq:kinematics}
\end{equation}

Из этих уравнений можно получить уравнение траектории, исключив параметр времени \(t\):

\begin{equation}
    z = z_0 + \frac{v_{0z}}{v_{0x}}(x - x_0) - \frac{g}{2 v_{0x}^2}(x - x_0)^2
    \label{eq:trajectory}
\end{equation}

Как видно из формулы \eqref{eq:trajectory}, траектория тела в однородном гравитационном поле представляет собой параболу.

\subsection{Учет сопротивления воздуха}

В более реалистичной постановке необходимо учитывать силу сопротивления воздуха. В большинстве практических случаев сила сопротивления пропорциональна квадрату скорости \cite{butenin2023}:

\begin{equation}
    \vec{F}_{\text{сопр}} = -k v \vec{v} = -c \rho S |\vec{v}| \vec{v}
    \label{eq:drag_force}
\end{equation}

где:
\begin{itemize}
    \item \(\rho\) — плотность воздуха;
    \item \(S\) — площадь поперечного сечения тела;
    \item \(c\) — коэффициент аэродинамического сопротивления;
    \item \(k = c\rho S\) — обобщенный коэффициент сопротивления.
\end{itemize}

Тогда система дифференциальных уравнений движения принимает вид:

\begin{equation}
\begin{cases}
    m\ddot{x} = -k \dot{x} \sqrt{\dot{x}^2 + \dot{y}^2 + \dot{z}^2} \\
    m\ddot{y} = -k \dot{y} \sqrt{\dot{x}^2 + \dot{y}^2 + \dot{z}^2} \\
    m\ddot{z} = -mg - k \dot{z} \sqrt{\dot{x}^2 + \dot{y}^2 + \dot{z}^2}
\end{cases}
\label{eq:drag_system}
\end{equation}

Данная система не имеет аналитического решения в элементарных функциях, что делает необходимым применение численных методов \cite{hairer2020}.